\documentclass[11pt]{article}

% ── Packages ─────────────────────────────────────────────────────────────────
\usepackage[margin=1in]{geometry}
\usepackage{amsmath,amssymb}
\usepackage{booktabs}
\usepackage{graphicx}
\usepackage{hyperref}
\usepackage{microtype}
\usepackage{listings}
\usepackage{xcolor}
\usepackage{caption}
\usepackage{subcaption}
\usepackage{parskip}
\usepackage{natbib}
\usepackage{enumitem}

% ── Hyperref setup ────────────────────────────────────────────────────────────
\hypersetup{
  colorlinks=true,
  linkcolor=blue!60!black,
  citecolor=blue!60!black,
  urlcolor=blue!60!black
}

% ── Code listing style ────────────────────────────────────────────────────────
\lstset{
  language=Python,
  basicstyle=\ttfamily\footnotesize,
  keywordstyle=\color{blue!70!black}\bfseries,
  stringstyle=\color{red!60!black},
  commentstyle=\color{gray},
  backgroundcolor=\color{gray!8},
  frame=single,
  framerule=0.4pt,
  rulecolor=\color{gray!40},
  breaklines=true,
  breakatwhitespace=true,
  columns=flexible,
  keepspaces=true,
  showstringspaces=false,
  xleftmargin=4pt,
  xrightmargin=4pt
}

% ── Title ─────────────────────────────────────────────────────────────────────
\title{%
  \textbf{Angular Manifold Routing: Sublinear Compute Reduction\\
  via Hopf-Base Sector Discretization}
}

\author{%
  Casey Allard\\[2pt]
  \textit{Independent Researcher}\\[2pt]
  \small\texttt{caseyallard@hotmail.com}
}

\date{March 2026}

% ─────────────────────────────────────────────────────────────────────────────
\begin{document}
\maketitle

% ── Abstract ─────────────────────────────────────────────────────────────────
\begin{abstract}
We show that the same geometric property of normalized token embeddings that enables
TurboQuant's extreme data compression~\citep{turboquant2026} also enables sublinear
routing computation in transformer-style architectures.
TurboQuant demonstrated that L2-normalized embeddings on the unit hypersphere are
\emph{angularly non-uniform}, and that random rotation destroys this structure for
near-optimal scalar quantization.
We independently developed the complementary direction: a fixed Hopf fibration map
exploits the same angular non-uniformity to concentrate token routing into a sparse
subset of available expert sectors.

Across 169 experimental increments on a PPMI-SVD semantic proxy, we establish:
(1)~the routing footprint scales as $K^{0.572}$ versus $K^{1.0}$ for dense routing;
(2)~this advantage grows with $K$ and persists at $K{=}5000$ (ratio $2.6$--$2.8\times$);
(3)~the mechanism is purely angular and norm-invariant ($\Delta\alpha < 0.015$ across
L1/L2/L3/L4 normalizations);
and (4)~a fixed 4D Hopf geometry outperforms 100D adaptive $K$-means on semantically
structured embeddings.

We validate the mechanism in a small trainable language model (PTB, confirmed 2 seeds,
4000 steps): fixed geometric routing replaces learned gating with only 8\% validation
perplexity cost and no learned gate matrix, while using 46 of 64 effective expert paths at
convergence ($1.4\times$ more efficient than dense routing).
A second-dataset replication on WikiText-2 (confirmed 2 seeds) finds a
HOPF/BASELINE ratio of $1.081$---numerically identical to the PTB confirmed ratio---under identical
training conditions.
A key honest finding: the specific Hopf sector geometry does not add advantage over
randomly-permuted fixed routing in trainable LM settings---experts co-adapt their
weights---but the 8\% PPL gap is stable across both datasets and both seed counts.
This result is scoped to the 2-layer toy-scale trainable setting and should not be
read as a claim of broad MoE replacement or large-scale transformer substitution.

We provide a standalone Python script (numpy only, $<60$~s) reproducing the core result.
Taken together with TurboQuant, this work suggests the angular non-uniformity of
embeddings has engineering consequences in both data compression and routing computation.
\end{abstract}

% ── 1. Background ─────────────────────────────────────────────────────────────
\section{Background: TurboQuant and Angular Non-Uniformity}

\subsection{The geometric fact}

\citet{turboquant2026} demonstrated a previously underappreciated property of modern
language model embeddings: L2-normalized token representations on the unit hypersphere
are \emph{angularly non-uniform}.
Projecting hidden states or token embeddings to the sphere reveals that they cluster
into preferred angular regions rather than distributing uniformly.

TurboQuant's engineering response: \textbf{destroy} this structure via random orthogonal
rotation, producing near-independent scalar components amenable to per-dimension scalar
quantization at extreme bit-depths.
This achieves state-of-the-art KV-cache and vector-index compression, validated at
production scale.

\subsection{Our parallel, complementary choice}

We independently arrived at the same geometric fact through a different path---not
compression, but routing efficiency.
If token embeddings are angularly non-uniform, a \emph{fixed sector map} that divides
the sphere into regions will receive non-uniform routing: some sectors capture many
semantically similar tokens while others stay empty.
This produces a sparse routing distribution, reducing the effective compute footprint.

Our engineering response: \textbf{exploit} the angular non-uniformity via a fixed Hopf
fibration map to create sparse, semantically coherent routing.

\begin{table}[h]
\centering
\caption{TurboQuant vs.\ this work: same geometric fact, opposite engineering choice.}
\label{tab:comparison}
\begin{tabular}{lll}
\toprule
                    & \textbf{TurboQuant} & \textbf{This work} \\
\midrule
Angular structure   & Destroyed (random rotation)   & Exploited (Hopf map)          \\
Purpose             & Data compression              & Compute reduction             \\
Mechanism           & Near-independent scalars      & Sparse sector routing         \\
Scale validated     & Production LLMs               & Proxy + 2-layer LM            \\
Claim               & Store/transmit fewer bits     & Compute fewer routing paths   \\
\bottomrule
\end{tabular}
\end{table}

\subsection{Independent parallel development}

Our research chain (INC-0143 through INC-0173) predates the TurboQuant public release.
TurboQuant provides independent corroboration of the geometric fact at production scale
and motivates framing our contribution as its complementary extension.

% ── 2. The Mechanism ──────────────────────────────────────────────────────────
\section{The Mechanism: Hopf Angular Routing}

\subsection{The winning algorithm}

After testing more than 20 routing variants---adaptive $K$-means, shell-aware variants,
product phase maps, complex transport, and SO(8) learned rotations---a single method
survived all falsification attempts: \texttt{sector\_mode = phase4d\_hopf\_transport},
$\lambda{=}1.0$, L2-normalization, single shell.

The algorithm assigns a routing sector to any L2-normalized input
$\mathbf{x} \in \mathbb{R}^D$ in four steps.

\textbf{Step 1 -- Project to 4D Hopf subspace:}
$a = x_{i},\; b = x_{j},\; c = x_{k},\; d = x_{l}$
(indices selected once on a held-out proxy; canonical values: $(3, 65, 2, 21)$).

\textbf{Step 2 -- Hopf base coordinates:}
\begin{align}
  \rho_1 &= \sqrt{a^2+b^2},\quad \rho_2 = \sqrt{c^2+d^2}, \nonumber\\
  \chi_u &= \frac{\rho_2^2}{\rho_1^2+\rho_2^2}\in[0,1], \quad
  \delta  = \arctan2(b,a)-\arctan2(d,c)\in[-\pi,\pi], \nonumber\\
  \alpha  &= \tfrac{1}{2}\bigl(\arctan2(b,a)+\arctan2(d,c)\bigr)\in[-\pi,\pi].
  \label{eq:hopf_coords}
\end{align}

\textbf{Step 3 -- Phase transport (Levi-Civita correction):}
\begin{equation}
  \chi = \arcsin\!\left(\frac{\rho_2}{\sqrt{\rho_1^2+\rho_2^2}}\right), \qquad
  \alpha_t = \alpha + \tfrac{1}{2}\lambda\cos(2\chi)\,\delta.
  \label{eq:transport}
\end{equation}

\textbf{Step 4 -- Assign sector} across budget $K$:
allocate bins $(k_\chi,k_\delta,k_\alpha)$ with product $\le K$; sector $= b_\chi \cdot k_\delta k_\alpha + b_\delta \cdot k_\alpha + b_{\alpha_t}$.

\subsection{Mechanism isolation (INC-0147)}

The raw fiber mean phase $\alpha = \frac{1}{2}(\theta_1+\theta_2)$ is the primary
mechanism (Table~\ref{tab:mechanism}).

\begin{table}[h]
\centering
\caption{Mechanism isolation on PPMI-SVD proxy.
Higher Gini ratio = more concentrated routing = better compute reduction.}
\label{tab:mechanism}
\begin{tabular}{lcc}
\toprule
Variant                                        & Gini ratio (rel\_diff) & Increment \\
\midrule
HOPF\_BASE ($\chi+\delta$ only, no $\alpha$)   & 46.0\%                 & INC-0147  \\
HOPF $\lambda{=}0$ (raw $\alpha$, no transport)& 66.7\% ($+$20.6 pp)    & INC-0147  \\
HOPF $\lambda{=}1.0$ (canonical)               & 68.7\% ($+$2 pp LC)    & INC-0147  \\
\bottomrule
\end{tabular}
\end{table}

The raw fiber phase $(\theta_1+\theta_2)/2$ provides nearly all the gain;
the Levi-Civita transport adds only ${\approx}2$ percentage points.

\subsection{Fixed geometry vs.\ adaptive (INC-0144)}

Fixed 4D Hopf geometry outperforms 100D adaptive $K$-means: rel\_diff 31.2\%
(stable, 5 seeds) vs.\ 3.1\% (variable, $-5.8\%$ to $+12.0\%$).
The Hopf map captures fiber geometry of the $S^3$ manifold; Euclidean cluster
centroids miss this structure.

% ── 3. Scaling Law ────────────────────────────────────────────────────────────
\section{Scaling Law and Hardware Consequences}

\subsection{The canonical scaling law}

On the PPMI-SVD proxy ($>$160 runs, 7 seeds, $K=25$ to $5000$):

\begin{table}[h]
\centering
\caption{Routing scaling laws. HOPF = canonical Hopf with full transport; PERM = structure-destroyed
control (same marginal statistics, randomly permuted sectors).}
\label{tab:scaling}
\begin{tabular}{llcc}
\toprule
Routing         & Scaling law                                & $R^2$  & Range \\
\midrule
HOPF            & $\hat{e} = 2.957 \cdot K^{\mathbf{0.572}}$ & 0.963  & $K$=25--400 \\
HOPF            & $\hat{e} \approx K^{0.657}$                 & 0.991  & $K$=25--5000 \\
PERM            & $\hat{e} = 1.664 \cdot K^{0.814}$           & 0.993  & $K$=25--5000 \\
DENSE           & $\hat{e} = K^{1.0}$                         & 1.000  & all $K$ \\
\bottomrule
\end{tabular}
\end{table}

Figure~\ref{fig:scaling} shows the log-log scaling curves.
The Hopf advantage grows with $K$: the PERM/ORIG compression ratio reaches
$2.21\times$ at $K{=}400$, $2.55\times$ at $K{=}1000$, and stabilises at
$2.6$--$2.8\times$ across $K{=}600$--$5000$ (INC-0170).

\begin{figure}[h]
\centering
\includegraphics[width=0.72\linewidth]{figures/fig_scaling_law.pdf}
\caption{Log-log routing footprint ($\mathrm{eff}_b$) vs.\ $K$.
  HOPF (blue circles) scales as $K^{0.572}$ vs.\ $K^{1.0}$ for dense routing.
  The PERM control (orange squares) scales at $K^{0.814}$---intermediate,
  confirming angular \emph{structure} (not mere non-uniformity) is the mechanism.
  Lines show parametric fits; points are confirmed experimental means.}
\label{fig:scaling}
\end{figure}

\subsection{Hardware proxy (INC-0165, 80 runs, 5 seeds)}

Simulating routing cache pressure:
HOPF achieves $3.0$--$4.9\times$ lower effective cost and $2.5$--$2.9\times$
fewer LRU-16 cache misses vs.\ dense routing.
Phase transport amplifies the hardware advantage by $1.5$--$2.4\times$ vs.\ BASE.
Gains do not saturate through $K{=}5000$.

\subsection{Norm invariance (INC-0168)}

The scaling exponent is norm-invariant: $\alpha{=}0.572$ varies by $<0.015$ across L1,
L2, L3, L4 normalizations.
The mechanism is purely angular.
Shell stratification (radial layers) degrades the Gini ratio by 19\% and is excluded.

% ── 4. Transformer Routing ────────────────────────────────────────────────────
\section{Extension to Transformer Routing}

\subsection{Experiment design}

Architecture: 2-layer transformer, $h{=}128$, 4 attention heads, vocab$=$5000,
seq\_len$=$32, $K{=}64$ experts.
Three FFN conditions (all $K{=}64$, $\ge 5.6$M parameters):

\begin{table}[h]
\centering
\caption{LM routing conditions.}
\label{tab:conditions}
\begin{tabular}{lll}
\toprule
Condition     & Gate             & Routing \\
\midrule
BASELINE      & Learned $K{\times}D$ + softmax $\to$ top-1 & Supervised \\
HOPF-ROUTED   & None             & L2-norm(hidden) $\to$ Hopf sector \\
PERMUTED      & None             & Hopf sectors randomly permuted \\
\bottomrule
\end{tabular}
\end{table}

Training: 4000 steps, cosine LR $3\times10^{-4}$, batch 64.
Confirmed: 2 seeds, PTB and WikiText-2.

\subsection{Results: PTB (INC-0171, 2 seeds confirmed)}

\begin{table}[h]
\centering
\caption{PTB LM results (2-seed mean, 4000 steps). eff$_b$ = effective expert paths
out of $K{=}64$. Lower eff$_b$ relative to $K$ = more concentrated routing.}
\label{tab:ptb}
\begin{tabular}{lcccc}
\toprule
Method        & Val PPL & eff$_b$ & Params    & eff\_ratio vs DENSE \\
\midrule
BASELINE      & \textbf{152.26} & 43.19 & 5,660,672 & $1.48\times$          \\
HOPF-ROUTED   & \textbf{164.54} & 45.96 & 5,644,288 & $\mathbf{1.39\times}$ \\
PERMUTED      & \textbf{164.41} & 45.81 & 5,644,288 & $1.40\times$          \\
\bottomrule
\end{tabular}
\end{table}

HOPF/BASELINE PPL ratio: \textbf{1.081} (2 seeds, both within $\le1.10$ threshold).\\
HOPF vs.\ PERMUTED: $\Delta\text{ppl} = +0.13$ (statistically indistinguishable).\\
eff\_ratio decline: $1.65\times$ at step 2000 $\to$ $1.39\times$ at convergence
(experts expand sector coverage during training; paper reports convergence value).

\subsection{Results: WikiText-2 (INC-0173, 2 seeds confirmed)}

\begin{table}[h]
\centering
\caption{WikiText-2 replication (2-seed detail). Identical setup to PTB.}
\label{tab:wt2}
\begin{tabular}{lccccc}
\toprule
Method & Seed 0 PPL & Seed 1 PPL & Mean PPL & Mean eff$_b$ \\
\midrule
BASELINE    & 106.96 & 104.74 & 105.85 & 35.37 \\
HOPF-ROUTED & 113.65 & 115.25 & 114.45 & 40.94 \\
PERMUTED    & 113.86 & 115.09 & 114.48 & 42.28 \\
\bottomrule
\end{tabular}
\end{table}

\begin{table}[h]
\centering
\caption{Cross-dataset summary. WT2 confirms the PTB result at equal 2-seed standard.}
\label{tab:crossdataset}
\begin{tabular}{lcc}
\toprule
Metric                      & PTB (INC-0171) & WT2 (INC-0173) \\
\midrule
HOPF/BASELINE PPL ratio     & 1.081          & \textbf{1.081} \\
HOPF vs.\ PERMUTED $|\Delta\text{ppl}|$ & 0.13  & \textbf{0.03}  \\
HOPF vs.\ DENSE eff\_ratio  & $1.39\times$   & $\mathbf{1.56\times}$ \\
BASELINE mean eff$_b$       & 44.00          & 35.37          \\
HOPF mean eff$_b$           & 45.96          & 40.94          \\
Seeds confirmed             & 2              & 2              \\
\bottomrule
\end{tabular}
\end{table}

\subsection{Interpretation}

\textbf{Finding 1 -- Fixed geometric routing replaces learned gating at 8\% PPL cost
(confirmed, cross-dataset).}
HOPF-ROUTED achieves within-8\% validation perplexity of the learned gating baseline
with no learned gate matrix, confirmed across 2 seeds and two datasets.
The HOPF/BASELINE ratio of 1.081 is numerically identical on PTB and WT2, an unusually tight cross-dataset replication.

\textbf{On the correct baseline.}
The BASELINE uses learned top-1 gating \emph{without} load-balance auxiliary loss.
A control experiment (INC-0172, killed on design grounds) showed that imposing
Switch-style~\citep{fedus2022switch} load-balance auxiliary loss drives routing to
near-uniform distribution (eff$_b{=}62.77$ vs.\ BASELINE eff$_b{=}44$) and slightly
degrades perplexity.
This establishes that \emph{native learned concentration}---not imposed balance---is the
correct architectural regime for this comparison.
HOPF and BASELINE both operate in the same concentration regime; the 8\% gap is the cost
of fixing the gate in that regime.

\textbf{Finding 2 -- Hopf geometry is not distinguished from random fixed routing in
trainable LM (confirmed).}
HOPF and PERMUTED are statistically indistinguishable across both datasets (PTB
$\Delta{=}0.13$~ppl; WT2 $|\Delta|{=}0.03$~ppl).
In a trainable LM, expert weight matrices co-adapt to any fixed routing scheme.
The proxy experiments (Sections~2--3) reveal genuine angular structure on static
embeddings; in end-to-end training, this structure advantage is absorbed.
The claim is precisely scoped to the efficiency transfer, not to geometry-specific
quality advantage.

% ── 5. Discussion ─────────────────────────────────────────────────────────────
\section{Discussion, Reproducibility, and Scope}

\subsection{Reproducing the core result}

\begin{lstlisting}[caption={Core routing function (numpy only, $<60$~s to reproduce Table~\ref{tab:scaling}).},
                   label={lst:routing}]
def hopf_transport_sectors(x, K, dims=(3, 65, 2, 21), lam=1.0):
    """Assign L2-normalised vectors to Hopf angular sectors."""
    a, b, c, d = x[:,dims[0]], x[:,dims[1]], x[:,dims[2]], x[:,dims[3]]
    rho1 = np.sqrt(a**2 + b**2)
    rho2 = np.sqrt(c**2 + d**2)
    chi_u = np.clip(rho2**2 / np.maximum(rho1**2 + rho2**2, 1e-30), 0, 1)
    chi   = np.arcsin(np.clip(rho2 / np.maximum(
                np.sqrt(rho1**2 + rho2**2), 1e-30), 0, 1))
    theta1, theta2 = np.arctan2(b, a), np.arctan2(d, c)
    delta = wrap_to_pi(theta1 - theta2)
    alpha = wrap_to_pi(0.5 * (theta1 + theta2))
    ta    = wrap_to_pi(alpha + 0.5 * lam * np.cos(2 * chi) * delta)
    kchi, kdelta, kalpha = allocate_triplet_bins_budget(K)
    bchi   = np.clip((chi_u * kchi).astype(int),   0, kchi   - 1)
    bdelta = np.clip(((delta  + np.pi)/(2*np.pi) * kdelta).astype(int),
                     0, kdelta - 1)
    balpha = np.clip(((ta + np.pi)/(2*np.pi) * kalpha).astype(int),
                     0, kalpha - 1)
    return bchi * kdelta * kalpha + bdelta * kalpha + balpha
\end{lstlisting}

Full reproduction:
\begin{verbatim}
python hopf_routing_demo.py   # numpy only; reproduces Table 3
\end{verbatim}

A standalone reproduction bundle, including \texttt{hopf\_routing\_demo.py},
figure-generation code, and manuscript source, is available at:
\url{https://doi.org/10.5281/zenodo.19243034}

\subsection{What this work does not prove}

We are explicit about the scope boundaries:

\begin{itemize}[leftmargin=1.4em,itemsep=2pt]

\item \textbf{Proxy scale only (Sections~2--3).}
The PPMI-SVD proxy captures semantic structure but uses an EMA update rule, not
backpropagation.
Large-scale LM validation of the proxy claims is future work;
TurboQuant provides partial independent corroboration of the geometric fact.

\item \textbf{Hopf geometry vs.\ random fixed routing not distinguished in LM training.}
The efficiency gain transfers; the specific sector geometry is absorbed by learning
(Finding~2, both datasets).

\item \textbf{Native concentration regime only.}
The BASELINE is top-1 gating without auxiliary loss.
The claim does not extend to Switch-Transformer-style MoE with explicit balance
constraints; the INC-0172 control confirmed that imposed balance yields near-uniform
routing, a different architectural regime.

\item \textbf{Attention routing not tested.}
All routing is in the FFN gate; Q/K angular routing is a separate problem.

\item \textbf{Single research group.}
No external replication for Sections~2--3.
TurboQuant is the closest independent corroboration of the underlying geometric fact.

\end{itemize}

\subsection{Broader perspective}

The convergence of TurboQuant and this work on the same geometric property---approached
from opposite engineering directions---invites a broader observation.
TurboQuant destroys angular non-uniformity to achieve compression;
this work exploits it to achieve routing sparsity;
the proxy experiments show it also governs semantic concentration and expert
specialization.
If the same angular geometry governs compression, concentration, and routing
substitution, then these may not be separate mechanisms at all, but partial
expressions of a deeper computational substrate: the way trained language models
organize information angularly on the hypersphere.
Mapping this substrate explicitly---rather than treating each application independently---
is the natural next research direction.

\subsection{Next steps}

The most natural extension is validating the efficiency transfer at larger model and
dataset scale.
More broadly, the convergence of TurboQuant and this work on a shared geometric property
invites investigation of whether the angular structure of normalized embeddings constitutes
a common substrate underlying multiple inference-time efficiency mechanisms;
we leave this to future work.

% ── References ────────────────────────────────────────────────────────────────
\bibliographystyle{plainnat}
\bibliography{references}

% ── Appendix ──────────────────────────────────────────────────────────────────
\appendix
\section{Evidence Chain Summary}

\begin{table}[h]
\centering
\caption{Key increments in the evidence chain. Full records are available in the released artifact bundle.}
\label{tab:evidence}
\small
\begin{tabular}{p{2.9cm}p{1.5cm}p{8.8cm}}
\toprule
Increment range     & Section & Finding \\
\midrule
INC-0143--0146      & \S2     & Fixed H\textsuperscript{4} Hopf outperforms 100D K-means;
                               Gini advantage established \\
INC-0147            & \S2     & Mechanism isolation: raw fiber $\alpha=(\theta_1+\theta_2)/2$
                               is primary mechanism (+20.6~pp) \\
INC-0148--0159      & \S3     & Training-time scaling, EMA update rule, $K{=}250$--$1000$ \\
INC-0160--0165      & \S3     & Hardware proxy: $3.0$--$4.9\times$ eff\_cost reduction,
                               $2.5$--$2.9\times$ LRU miss reduction \\
INC-0166--0168      & \S3     & Norm invariance ($\Delta\alpha < 0.015$), shell hierarchy
                               excluded, canonical law \\
INC-0169            & \S3     & Law freeze: $\hat{e} = 2.957 \cdot K^{0.572}$ (canonical HOPF) \\
INC-0170            & \S3     & Large-$K$: ratio $2.6$--$2.8\times$ at $K{=}600$--$5000$
                               (does not collapse) \\
INC-0171            & \S4     & LM integration (PTB, 2 seeds): HOPF within 8\% PPL,
                               eff$_b{=}46/64$, $1.39\times$ vs DENSE \\
INC-0172 (killed)   & \S4     & Imposed load-balance aux loss $\to$ near-uniform routing
                               (eff$_b{=}63$); confirms native concentration
                               as correct comparator \\
INC-0173            & \S4     & WT2 replication (2 seeds): ratio $1.081$ (mean),
                               identical to PTB; HOPF${\approx}$PERM ($|\Delta|{=}0.03$~ppl) \\
\bottomrule
\end{tabular}
\end{table}

\end{document}
